\chapter{Requirements Engineering}

This chapter establishes \emph{what} the AI-Driven Smart Home Energy Optimizer
(SHEO) must do, deliberately kept separate from \emph{how} it is built. Applying
established requirements-analysis discipline, it first walks through
the requirements engineering process, then reconstructs the client elicitation
session from which the requirements were drawn, and finally records the catalogue of
functional and non-functional requirements, the use-case model, the use-case
specifications, and the behavioural and data models. The artefacts produced here form
the agreed basis against which the system is later designed, constructed, and
acceptance-tested.

\section{The Requirements Engineering Process}
\label{sec:req-process}

Requirements engineering proceeds as a five-step pipeline whose output is a
Software Requirements Specification (SRS)~\cite{ieee29148}. SHEO followed this pipeline faithfully:

\begin{enumerate}
  \item \textbf{Feasibility study.} Before committing, the team checked that the
  product was technically, operationally, and economically achievable within the
  project timeline. The decisive constraint --- no physical Internet-of-Things (IoT)
  hardware --- was resolved by deciding to substitute a deterministic mock device
  simulator, which keeps the whole product demonstrable end-to-end while remaining
  hardware-ready. The conclusion was that a small, fully working system was feasible
  and preferable to a large half-built one.
  \item \textbf{Requirements collection (elicitation).} Requirements were gathered
  directly from the client in an offline discussion (Section~\ref{sec:elicitation}).
  The principal elicitation techniques applied were the \emph{structured interview}
  (asking the client to describe the devices, the desired controls, and the business
  priority) and \emph{observation / scenario walk-through} (reasoning about how a
  resident would inspect consumption, control a device, and act on a suggestion). A
  structured client-discussion approach --- posing sharp scoping questions to the
  client --- supplied the model for this step.
  \item \textbf{Analysis.} The raw notes were decoded, ambiguities were surfaced, and
  conflicts and gaps were negotiated into testable statements. For example, the
  client's loose phrase ``find habit and apply that rule'' was analysed into two
  separable obligations: \emph{mine} a habit and \emph{recommend} a rule, with the
  user retaining explicit control over whether it is ever applied.
  \item \textbf{Modelling.} The analysed requirements were captured in standard
  models: a use-case diagram and use-case specification tables for the functional
  view (Sections~\ref{sec:usecase}--\ref{sec:uc-spec}), and an
  entity-relationship diagram together with finite-state machines for the data and
  behavioural views (Section~\ref{sec:models}).
  \item \textbf{Definition (SRS).} Finally the requirements were written up as a
  structured catalogue of functional requirements (\req{REQ-4.x}, \req{REQ-6.1}),
  non-functional requirements (\req{NFR-*}), and business rules (\req{BR-1}--\req{BR-5}),
  each given a unique identifier so that it can be traced forward into design, code,
  and test. Standard SRS quality criteria --- that it be \emph{clear, accurate,
  relevant,} and \emph{complete} --- guided this write-up, and the requirement
  identifiers used throughout this chapter are exactly those carried into the
  traceability matrix.
\end{enumerate}

A central rule underpins the entire chapter: requirements
(``what'') are kept strictly separate from design (``how''). The statements below
describe the system's externally observable obligations; the architecture and source
modules that satisfy them are the subject of later chapters.

\section{Client Elicitation Session}
\label{sec:elicitation}

The requirements originate in an offline discussion with the project's client, an
IoT contractor. Reconstructed and decoded, the session ran as follows.

\textbf{Who the client is.} The client is the party
that \emph{commissions and funds} the build. Here that is an IoT contractor who
manufactures a range of AI capabilities and a range of IoT devices, and who hires the
team to deliver an application on top of them. Crucially, the contractor is a
\emph{business stakeholder, not a system actor}: it pays for and supplies the system
but does not operate the running software day to day. The parties who do operate it ---
the building management (the \emph{customer}) and the residents (the \emph{end-users})
--- are distinguished from the client throughout.

\textbf{What the client asked for.} The contractor ``has all kinds of IoT devices''
and ``wants to provide an application.'' Each device type exposes its own controls:
a fan can change speed and set a mode; an air-conditioner exposes a target temperature
and an on/off switch. The client did not want the application to hard-code one screen
per device. Instead it asked for an \emph{application programming interface that
decides which feature each device exposes}, so that the appropriate controls for any
given device are determined dynamically. Because no real hardware was available for
the build, the client agreed that the team should ``create some kind of a mock
device'' standing in for the physical ones.

\textbf{Automation and recommendations.} The client described automation in terms of
\emph{preconditions and postconditions} attached to a device's behaviour: a user
\emph{may} set these conditions if they wish, but --- in the client's words ---
``if they don't,'' the system should ``find a habit and apply that rule.'' Analysis
turned this into the habit-recommendation requirement: when the user has not authored
a rule, the system mines the resident's usage pattern and \emph{recommends} an
automation rule, which the user can then read, accept, or dismiss.

\textbf{The high-priority requirement.} The client stated one priority explicitly and
emphatically: \emph{energy saving}. The contractor wants the customer ``to see the
result of the actual electricity bill they saved when using the suggested rule.''
This single sentence fixes the product's centre of gravity: every suggested rule must
carry a concrete, monetary estimate of the saving it would achieve, expressed in
Vietnamese đồng (VND), and shown \emph{before} the rule is accepted. The estimation
must therefore be transparent and explainable rather than a black box --- a constraint
that shapes the whole optimization design.

\textbf{Clarifying questions.} Having received a terse client brief, the team returned
to the client with sharp questions that pinned down scope, roles, and edge cases before
designing anything. These clarifying questions resolved exactly the points the
client's notes left open: \emph{scope} (mock devices instead of real hardware, an
explainable rule engine instead of opaque machine learning), \emph{roles} (which party
is client, customer, and end-user, and which of these becomes a login actor), and
\emph{edge cases} (what happens when a device is offline, when a command is out of
range, or when a safety-critical device would be powered off automatically). Those
answers are precisely the obligations catalogued below.

\section{Functional Requirements}
\label{sec:fr}

The functional requirements are organised into six families, each tracing back to a
point raised in the elicitation session. Table~\ref{tab:fr} lists the most important
sub-requirements; every one carries a stable identifier (\req{REQ-4.x.y},
\req{REQ-6.1}) and a priority (H = high, M = medium) used in the traceability matrix.

{\small
\begin{xltabular}{\linewidth}{@{}l l Y c@{}}
\caption{Functional requirements catalogue (the principal sub-requirements).}\label{tab:fr}\\
\toprule
\textbf{ID} & \textbf{Family} & \textbf{Requirement (the system shall\,\ldots)} & \textbf{Pri.}\\
\midrule
\endfirsthead
\toprule
\textbf{ID} & \textbf{Family} & \textbf{Requirement (the system shall\,\ldots)} & \textbf{Pri.}\\
\midrule
\endhead
\multicolumn{4}{@{}l}{\textit{REQ-4.1 --- Real-time monitoring}}\\
\req{REQ-4.1.1} & Monitoring & display instantaneous power per device, refreshed within 5\,s. & H\\
\req{REQ-4.1.2} & Monitoring & report cumulative energy in kWh for today, the billing cycle, and a chosen range. & H\\
\req{REQ-4.1.3} & Monitoring & retain one-minute telemetry aggregates for at least 12 months. & H\\
\req{REQ-4.1.4} & Monitoring & mark a device unreachable after more than 60\,s of silence and exclude it from the live total. & H\\
\req{REQ-4.1.5} & Monitoring & rank the top consumers by day, week, and month. & H\\
\midrule
\multicolumn{4}{@{}l}{\textit{REQ-4.2 --- Capability-driven control}}\\
\req{REQ-4.2.1} & Control & expose, per device, a capability description that the front-end fetches \emph{before} rendering controls. & H\\
\req{REQ-4.2.2} & Control & provide on/off and per-type controls for plug, bulb, fan, and air-conditioner. & H\\
\req{REQ-4.2.3} & Control & validate each command against the device's capability range and reject anything out of range. & H\\
\req{REQ-4.2.4} & Control & return a command outcome --- \textsc{success}, \textsc{rejected}, or \textsc{timeout} --- within 5\,s. & H\\
\req{REQ-4.2.5} & Control & provide a mock simulator for plug, bulb, fan, air-conditioner, and sensor device types. & H\\
\midrule
\multicolumn{4}{@{}l}{\textit{REQ-4.3 --- Explainable rules \& auto-actions}}\\
\req{REQ-4.3.1} & Rules & support conditions over time, day, tariff, device state, and sensor readings. & H\\
\req{REQ-4.3.2} & Rules & require a rule's action to match the target device's capability. & H\\
\req{REQ-4.3.3} & Rules & detect conflicts between enabled rules whose time windows overlap. & H\\
\req{REQ-4.3.4} & Rules & let the user enable, disable, edit, and delete rules. & H\\
\req{REQ-4.3.5} & Rules & log every rule execution in an immutable audit record. & H\\
\req{REQ-4.3.6} & Rules & keep automatic action opt-in (off by default) and offer a 2-minute undo. & M\\
\midrule
\multicolumn{4}{@{}l}{\textit{REQ-4.4 --- Habit recommendations}}\\
\req{REQ-4.4.1} & Recommend & generate recommendations only after at least 7 days of telemetry. & H\\
\req{REQ-4.4.2} & Recommend & express each recommendation as a readable \texttt{WHEN--THEN} rule with a rationale and data window. & H\\
\req{REQ-4.4.3} & Recommend & rank recommendations by VND saving and keep at most five active. & H\\
\req{REQ-4.4.4} & Recommend & convert a recommendation into a real rule only on explicit acceptance. & H\\
\req{REQ-4.4.5} & Recommend & suppress a dismissed recommendation for at least 30 days. & H\\
\midrule
\multicolumn{4}{@{}l}{\textit{REQ-4.5 --- VND bill-saving estimation (the client's high priority)}}\\
\req{REQ-4.5.1} & Savings & compute a 14-day hourly baseline energy profile per device. & H\\
\req{REQ-4.5.2} & Savings & estimate a saving as $\Sigma\,((\text{baseline}-\text{expected})\times\text{tariff})$. & H\\
\req{REQ-4.5.3} & Savings & show the VND estimate \emph{before} the user saves a rule or accepts a recommendation. & H\\
\req{REQ-4.5.4} & Savings & display the saving actually accrued so far in the current billing cycle (VND). & H\\
\req{REQ-4.5.5} & Savings & flag a rule for recalculation when measured saving drifts beyond $\pm$20\% of the estimate. & M\\
\midrule
\multicolumn{4}{@{}l}{\textit{REQ-6.1 --- Reporting}}\\
\req{REQ-6.1} & Reporting & export readings, rules, and savings as CSV. & M\\
\bottomrule
\end{xltabular}
}

Two features deserve emphasis because they encode the client's explicit wishes.
First, \req{REQ-4.2.1} captures the requirement for an interface ``deciding which
feature'' a device exposes: controls are rendered from a per-type capability
description rather than hard-coded, so a new device type is supported without new
control code. Second, \req{REQ-4.5.3} encodes the high-priority demand that the
customer \emph{see} the saving of a suggested rule before committing to it --- the
estimate must be visible at decision time, not merely after the fact.

\section{Non-Functional Requirements and Business Rules}
\label{sec:nfr}

Non-functional requirements are classified as \emph{product},
\emph{organisational}, or \emph{external} requirements. SHEO's product NFRs cover
performance (\req{NFR-PER}), safety (\req{NFR-SAF}), and security (\req{NFR-SEC}); the
external requirements are the deployment-level transport-security obligations; and the
business rules (\req{BR-1}--\req{BR-5}) are organisational constraints that the
business places on the system's behaviour. Tables~\ref{tab:nfr} and~\ref{tab:br}
record them.

\begin{table}[htbp]
\centering
\caption{Non-functional requirements (product / external).}
\label{tab:nfr}
\small
\begin{tabularx}{\linewidth}{@{}l l Y c@{}}
\toprule
\textbf{ID} & \textbf{Class} & \textbf{Requirement} & \textbf{Pri.}\\
\midrule
\req{NFR-PER-1} & Product (performance) & The dashboard shall refresh within 5\,s for up to 30 devices. & H\\
\req{NFR-PER-2} & Product (performance) & A control command shall return user-visible feedback within 2\,s (95th percentile). & H\\
\req{NFR-PER-3} & Product (performance) & The system shall sustain at least 100 concurrent simulated homes on one node. & H\\
\req{NFR-SAF-1} & Product (safety) & At most one air-conditioner compressor command shall be issued per 3 minutes. & H\\
\req{NFR-SAF-2} & Product (safety) & The system shall never automatically power off a safety-critical device. & H\\
\req{NFR-SAF-3} & Product (safety) & Authoring shall warn on temperature, safety-critical, and unattended-off actions. & H\\
\req{NFR-SEC-1} & External (deployment) & All traffic shall use TLS 1.2+ (HTTPS / WSS). & H\\
\req{NFR-SEC-2} & Product (security) & The system shall enforce authentication and role-based access control. & H\\
\req{NFR-SEC-3} & Product (security) & No password or token shall be stored in plaintext. & H\\
\req{NFR-SEC-4} & Product (security) & Access to readings and logs shall be scoped to the user's own unit; the Administrator may view every unit in the building. & H\\
\req{NFR-SEC-5} & Product (security) & The building overview shall disclose only aggregate per-unit metrics (live load, energy, bill) and device \emph{reachability}, never a resident's per-device on/off state --- data minimisation, so oversight does not become behavioural surveillance. & M\\
\bottomrule
\end{tabularx}
\end{table}

\begin{table}[htbp]
\centering
\caption{Business rules --- organisational constraints (\req{BR-1}--\req{BR-5}).}
\label{tab:br}
\small
\begin{tabularx}{\linewidth}{@{}l Y c@{}}
\toprule
\textbf{ID} & \textbf{Business rule} & \textbf{Pri.}\\
\midrule
\req{BR-1} & Only the Administrator (building owner) may onboard a resident by selling them a unit (provisioned with its default device package) or offboard a resident, after which the unit becomes vacant while its devices and billing history are retained; a Resident cannot add or remove devices. & H\\
\req{BR-2} & A recommendation remains advisory and has no effect until it is explicitly accepted. & H\\
\req{BR-3} & Automatic action is disabled by default; the user must opt in per rule. & H\\
\req{BR-4} & Deleting a rule shall preserve its execution history. & H\\
\req{BR-5} & The tariff is manually configurable and defaults to VND. & M\\
\bottomrule
\end{tabularx}
\end{table}

The safety NFRs and \req{BR-2}/\req{BR-3} together express a deliberate caution
principle that runs through the requirements: the system advises and estimates, but it
acts on a unit's devices only when a human has explicitly authorised it, and
never in a way that could compromise a safety-critical device.

\section{Use-Case Model}
\label{sec:usecase}

The functional requirements above are organised into a use-case model that views the
system as a single black box used by a small set of external actors. By convention,
each actor is named with a \emph{singular role noun} (not an
individual or a job title), and each use case is named \emph{verb + noun}.

\fig{fig04-use-case.png}{Use-case model. The deployment models one apartment building. The Administrator (building owner) monitors a building overview across every unit, configures the building-wide tariff, and onboards or offboards residents by selling a unit or returning it to vacant --- but never manually operates a unit's devices and authors no automation rules; the Resident operates only their own unit's devices and may author rules and accept recommendations for that unit; the Developer/Tester exercises mock devices on a maintenance unit; the System Scheduler autonomously fires rules and mines habits. The Client (IoT contractor) and Customer (building management) are business stakeholders outside the boundary --- the Customer operates the system \emph{as} the Administrator. Accepting a recommendation $\langle\langle$include$\rangle\rangle$ authoring a rule, which in turn $\langle\langle$include$\rangle\rangle$ command validation and a VND estimate.}{fig:usecase}

The model has four actors:

\begin{itemize}
  \item \textbf{Administrator} --- the building owner, realised by the customer (building
  management). Owns no unit and never operates individual devices: monitors a read-only
  building overview across all units (resident roster, per-unit usage, building totals),
  configures the building-wide tariff, onboards residents by selling them a unit
  (provisioned with its default device package), and offboards a resident --- a soft
  removal that returns the unit to vacant while retaining its devices and history.
  \item \textbf{Resident} --- the primary end-user (a tenant of one unit). Views their own
  unit's consumption and savings, controls every operable device in that unit, and may
  author automation rules and accept or dismiss recommendations for it; no other unit is
  visible to them.
  \item \textbf{Developer/Tester} --- a maintenance role that may operate any device
  for diagnostics, exercises the mock devices, and uses the offline-test hook to
  prepare and reproduce demos.
  \item \textbf{System Scheduler} --- a non-human (time) actor that autonomously fires
  due rules and runs the periodic habit-mining and savings-accrual jobs.
\end{itemize}

The model uses the $\langle\langle$include$\rangle\rangle$ relationship for mandatory
sub-behaviour: \emph{Review \& accept recommendation} includes \emph{Author automation
rule} (an accepted recommendation becomes a rule), which itself includes the common
command-validation and VND-estimation steps. The Client and Customer
are drawn \emph{outside} the system boundary as business stakeholders, since neither
the funding contractor nor the adopting management is itself a login actor --- the
customer interacts with the running system only in its Administrator role.

\section{Use-Case Specifications}
\label{sec:uc-spec}

The principal use cases are detailed in specification tables giving
preconditions, the main flow, alternate and error-handling flows, and postconditions.
Three core use cases are specified below: controlling a device (UC-02), reviewing and
accepting a recommendation (UC-05), and authoring an automation rule (UC-04).

\begin{table}[htbp]
\centering
\caption{Use-case specification --- UC-02 Control device.}
\label{tab:uc02}
\small
\begin{tabularx}{\linewidth}{@{}L{0.20\linewidth} Y@{}}
\toprule
\textbf{Field} & \textbf{Description}\\
\midrule
UC code & UC-02\\
Name & Control device\\
Primary actor & Resident\\
Secondary actor & --- (the targeted device; the simulator stands in for hardware)\\
Preconditions & The actor is authenticated; the device exists in the actor's own unit; the actor holds control permission for it.\\
Main flow &
  \begin{enumerate}
    \item The actor opens the device's detail view.
    \item The system fetches the device's capability description (\req{REQ-4.2.1}).
    \item The system renders the controls generically from each control's kind (toggle / range / enumeration).
    \item The actor issues a command (e.g.\ set fan speed, set air-conditioner target temperature).
    \item The system validates the command against the capability range and the safety guard (\req{REQ-4.2.3}).
    \item The system applies the state change and records telemetry.
    \item The system returns the outcome \textsc{success} with the new state within 5\,s (\req{REQ-4.2.4}).
  \end{enumerate}\\
Alternate flow &
  \textit{A1 --- Out-of-range or schema-mismatched command:} at step 5 validation
  fails; the system makes no state change and returns \textsc{rejected} with a clear
  message; flow ends. \textit{A2 --- Safety guard:} a safety-breaching command (e.g.\
  exceeding the air-conditioner compressor rate limit) is returned as
  \textsc{rejected} / \textsc{skipped} and nothing changes (\req{NFR-SAF-1}).\\
Error-handling flow &
  \textit{E1 --- Device offline:} if the device has been silent beyond the timeout,
  the system returns \textsc{timeout}, changes nothing, and surfaces an offline
  indication (\req{REQ-4.2.4}, \req{REQ-4.1.4}).\\
Postcondition &
  On \textsc{success} the device is in the commanded state and a reading is logged; on
  \textsc{rejected}/\textsc{timeout} the device state is unchanged.\\
\bottomrule
\end{tabularx}
\end{table}

\begin{table}[htbp]
\centering
\caption{Use-case specification --- UC-05 Review \& accept recommendation.}
\label{tab:uc05}
\small
\begin{tabularx}{\linewidth}{@{}L{0.20\linewidth} Y@{}}
\toprule
\textbf{Field} & \textbf{Description}\\
\midrule
UC code & UC-05\\
Name & Review \& accept recommendation\\
Primary actor & Resident\\
Secondary actor & System Scheduler (mines the habits that produce recommendations)\\
Preconditions & The actor is authenticated; at least one active recommendation exists, derived from $\ge$\,7 days of telemetry (\req{REQ-4.4.1}).\\
Main flow &
  \begin{enumerate}
    \item The actor opens the Recommendations view.
    \item The system lists active recommendations, ranked by VND saving, at most five (\req{REQ-4.4.3}).
    \item For each, the system shows the readable \texttt{WHEN--THEN} rule, a rationale, the data window, and the monthly VND saving (\req{REQ-4.4.2}).
    \item The actor selects \emph{Accept}.
    \item The system converts the recommendation into a real automation rule ($\langle\langle$include$\rangle\rangle$ UC-04) and marks the recommendation accepted (\req{REQ-4.4.4}).
  \end{enumerate}\\
Alternate flow &
  \textit{A1 --- Dismiss:} at step 4 the actor selects \emph{Dismiss}; the system
  suppresses the same recommendation signature for at least 30 days and creates no
  rule (\req{REQ-4.4.5}). \textit{A2 --- Cap exceeded:} when more than five
  recommendations would be active, the lowest-saving surplus is expired to honour the
  cap.\\
Error-handling flow &
  \textit{E1 --- Insufficient data:} if fewer than 7 days of telemetry exist for a
  device, no recommendation is offered for it and the view explains why.\\
Postcondition &
  On accept, a new automation rule exists (auto-action off by default) and the
  recommendation is accepted; on dismiss, the recommendation is suppressed and no rule
  is created.\\
\bottomrule
\end{tabularx}
\end{table}

\begin{table}[htbp]
\centering
\caption{Use-case specification --- UC-04 Author automation rule.}
\label{tab:uc04}
\small
\begin{tabularx}{\linewidth}{@{}L{0.20\linewidth} Y@{}}
\toprule
\textbf{Field} & \textbf{Description}\\
\midrule
UC code & UC-04\\
Name & Author automation rule\\
Primary actor & Resident\\
Secondary actor & --- (the Optimization service supplies the VND estimate)\\
Preconditions & The actor is authenticated and holds control permission for the target device.\\
Main flow &
  \begin{enumerate}
    \item The actor opens the rule editor and chooses a target device.
    \item The actor specifies the \textsc{when} condition (time, day, tariff, state, or sensor) and the \textsc{then} action, optionally with an \textsc{until} stop condition (\req{REQ-4.3.1}).
    \item The system validates the action against the device capability (\req{REQ-4.3.2}) and checks for conflicts with overlapping enabled rules (\req{REQ-4.3.3}).
    \item The system computes and displays the VND saving estimate \emph{before} the rule is saved (\req{REQ-4.5.3}).
    \item The actor leaves auto-action off (default) or opts in, then saves.
    \item The system persists the rule in the Enabled state.
  \end{enumerate}\\
Alternate flow &
  \textit{A1 --- Conflict detected:} at step 3 the system reports the conflicting rule
  and the time overlap; the rule is not silently saved until the actor resolves it.
  \textit{A2 --- Opt-in to auto-action:} the actor enables automatic application; the
  system records the opt-in and later fires the rule subject to the safety guard.\\
Error-handling flow &
  \textit{E1 --- Action exceeds capability:} a \textsc{then} action outside the
  device's capability is rejected at validation with a clear message; the rule is not
  saved.\\
Postcondition &
  A valid rule is stored (Enabled), carrying its VND estimate; on validation failure no
  rule is created.\\
\bottomrule
\end{tabularx}
\end{table}

\section{Behavioural and Data Models}
\label{sec:models}

The modelling step of the requirements process also produces a data model and a set
of behavioural models, which together fix the structure of the system's state and the
lifecycles of its key objects. These are requirements-level models; their realisation
in the database schema and the running services is treated in later chapters.

\subsection{Data model}

Figure~\ref{fig:erd} gives the entity-relationship model. The deployment is one apartment
building; each \emph{home} is a \emph{unit} occupied by one Resident, and almost every
entity is scoped to a unit, which is the basis of the unit-scoped access control demanded
by \req{NFR-SEC-4}: a resident can never read another unit's readings or logs, while the
Administrator (building owner, with no unit of their own) may view every unit. A tariff
with no unit is the building-wide tariff. The model reifies the many-to-many control grant
between users and devices as a device-permission entity (the basis of \req{BR-1}), and it
deliberately keeps a rule's execution log independent of the rule so that deleting a rule
preserves its history, as \req{BR-4} requires.

\fig{fig08-erd.png}{Entity-relationship diagram. The deployment models one building; each Home is a unit, and almost every entity is scoped to a unit (the basis of unit-scoped access control, \req{NFR-SEC-4}). The Administrator has no unit (nullable \texttt{home\_id}) and oversees all units; a tariff with no unit is building-wide. The device-permission entity reifies the user--device control grant, and the rule-to-execution link carries no cascade so the execution log survives rule deletion (\req{BR-4}).}{fig:erd}

\subsection{Behavioural models}

Three finite-state machines specify the lifecycles of the system's most important
objects, modelling the requirements as observable state transitions.

\fig{fig12-state-rule.png}{State machine --- rule lifecycle. A rule is authored as a Draft, becomes Enabled or Disabled, transits a transient Fired state (auto-applied firings may be Undone within the 2-minute window), enters NeedsRecalc on $\pm$20\% drift, and on deletion is terminal yet preserves its execution history (\req{REQ-4.3}, \req{BR-4}).}{fig:state-rule}

The rule lifecycle (Figure~\ref{fig:state-rule}) captures \req{REQ-4.3.4} and
\req{REQ-4.3.6}: a rule moves between Enabled and Disabled, is fired by the scheduler,
and an auto-applied firing can be undone within two minutes.

\fig{fig13-state-recommendation.png}{State machine --- recommendation lifecycle. A mined recommendation becomes Active, then is Accepted (becoming a rule), Dismissed (suppressed for 30 days), or Expired when it falls outside the top-five cap (\req{REQ-4.4}).}{fig:state-rec}

The recommendation lifecycle (Figure~\ref{fig:state-rec}) models \req{REQ-4.4.3}--%
\req{REQ-4.4.5}: a recommendation is advisory while Active and leaves that state only
on an explicit user decision, becoming a rule on acceptance or suppressed for 30 days
on dismissal.

\fig{fig14-state-connectivity.png}{State machine --- device connectivity. A device flips to Offline after more than 60\,s of silence (or a forced offline) and back to Online when telemetry resumes; while offline it is excluded from the live home total (\req{REQ-4.1.4}).}{fig:state-conn}

The connectivity lifecycle (Figure~\ref{fig:state-conn}) realises \req{REQ-4.1.4}: a
silent device becomes Offline and is excluded from the live total until telemetry
resumes. A complementary activity-style model of the end-to-end user journey ---
signing in, then choosing to control a device, review a recommendation, or author a
rule --- is presented with the user-interface design in the later chapter, since it is
most naturally read alongside the screens it traverses.
